\section*{अध्याय \ref{ch:b2:magmatic-series} --- मैग्मीय श्रेणियाँ और विवर्तनिक परिवेश}

\begin{solution}{exo:b2:magmatic-series:1}
(क) $\ce{SiO2} = 49.60$, क्षार $2.22 + 0.51 = 2.73$: 45 और 52 के बीच और 5 से नीचे,
\emph{बेसाल्ट} (B)। (ख) 58.5 और $3.9 + 1.6 = 5.5$: 57 और 63 के बीच; वहाँ ऊपरी
सीमा (57, 5.9) से (63, 7.0) तक जाती है, और 58.5 पर $5.9 + 1.1\times1.5/6 = 6.18$ पर है;
5.5 उसके नीचे है: \emph{एंडेसाइट} (O2)। (ग) 72.0 और $4.0 + 4.2 = 8.2$:
69 पर ऊर्ध्वाधर रेखा से आगे, 8 से ऊपर और (69, 8) से (77, 0) तक की रेखा से ऊपर, जो
72 पर 5.0 पर है: \emph{रायोलाइट} (R)। (घ) 47.0 और $4.5 + 2.5 = 7.0$: 47 पर क्षेत्र U1
(45, 5) से (49.4, 7.3) तक की रेखा, जो $5 + 2.3\times2/4.4 = 6.05$ पर है, और
(45, 9.4) से (49.4, 7.3) तक की रेखा, जो $9.4 - 2.1\times2/4.4 = 8.45$ पर है, के बीच स्थित है। चट्टान
U1 में है: यदि उसके नॉर्म में \qty{10}{\percent} से कम ओलिवीन हो तो \emph{टेफ़्राइट},
अधिक हो तो \emph{बेसानाइट}।
\end{solution}

\begin{solution}{exo:b2:magmatic-series:2}
(क) ऑक्साइडों का योग 98.40 है और वाष्पशीलों सहित योग 100.00;
\ce{H2O} $1.40 < 2$ और \ce{CO2} $0.20 < 0.5$: विश्लेषण स्वीकार्य है। (ख) गुणक
$100/98.40 = 1.01626$: \ce{SiO2} 52.44, \ce{TiO2} 1.22, \ce{Al2O3} 16.36,
\ce{FeO_T} 9.96, \ce{MgO} 6.71, \ce{CaO} 8.33, \ce{Na2O} 3.56, \ce{K2O} 1.42
(योग 100.00); क्षार $3.56 + 1.42 = 4.98$। (ग) पहले: 51.60 और 4.90, 52 पर
बेसाल्ट--बेसाल्टी एंडेसाइट सीमा से ठीक नीचे: \emph{बेसाल्ट}। बाद में: 52.44 और 4.98;
52.44 पर O1 की ऊपरी सीमा, जो (52, 5) से (57, 5.9) तक जाती है, $5 + 0.9\times0.44/5
= 5.08$ पर है: \emph{बेसाल्टी एंडेसाइट} (O1)। पुनर्परिकलन ने चट्टान को एक सीमा के पार
खिसका दिया: किसी रेखा के निकट के नाम उस पर निर्भर करते हैं। (घ) 52.44 पर निम्न-K रेखा $0.3 +
0.045\times4.44 = 0.50$ पर है और उच्च-K रेखा $1.2 + 0.095\times4.44 = 1.62$ पर; 1.42
बीच में है: \emph{मध्यम-K}।
\end{solution}

\begin{solution}{exo:b2:magmatic-series:3}
(क) $Q + A + P = 90$: $Q = 25/90 = \qty{27.8}{\percent}$ और $100\,P/(A + P) =
100\times35/65 = 53.8$: क्षेत्र 3b, एक \emph{मॉन्ज़ोग्रेनाइट}, यानी ग्रेनाइट। बायोटाइट और
हॉर्नब्लेंड QAPF योग में शामिल नहीं होते। (ख) $Q + A + P = 65$: $Q = 5/65 =
\qty{7.7}{\percent}$ (5 और 20 के बीच) और $100\times55/60 = 91.7$ (90 से ऊपर): क्षेत्र
10*, क्वार्ट्ज़ डायोराइट, क्वार्ट्ज़ गैब्रो या क्वार्ट्ज़ ऐनॉर्थोसाइट। अकेला \omterm{def:b2:magmatic-series:qapf}{QAPF आरेख}
इनमें से चुनाव नहीं करता; यह प्लैजियोक्लेज़ और मैफ़िक खनिज (यहाँ
\qty{35}{\percent} क्लाइनोपाइरॉक्सीन) करते हैं।
\end{solution}

\begin{solution}{exo:b2:magmatic-series:4}
निम्न-K रेखा $0.3 + 0.045(\ce{SiO2} - 48)$ है, उच्च-K रेखा $1.2 +
0.095(\ce{SiO2} - 48)$। (क) 50 पर: 0.39 और 1.39; $0.40$ निम्न रेखा से ठीक ऊपर है:
\emph{मध्यम-K}। (ख) 56 पर: 0.66 और 1.96: \emph{मध्यम-K}। (ग) 62 पर: 0.93 और 2.53;
$2.90 > 2.53$: \emph{उच्च-K}। (घ) 66 पर: 1.11 और 2.91; $1.00 < 1.11$: \emph{निम्न-K}।
\end{solution}

\begin{solution}{exo:b2:magmatic-series:5}
(क) बेसाल्ट: $A + F + M = 3.2 + 14.5 + 4.5 = 22.2$; $A = 14.4$, $F = 65.3$, $M = 20.3$।
एंडेसाइट: $5.05 + 6.1 + 3.8 = 14.95$; $A = 33.8$, $F = 40.8$, $M = 25.4$। (ख) रेखा
कभी $F = 54.5$ से आगे नहीं जाती; बेसाल्ट, जिसका $F = 65.3$ है, F की ओर पड़ता है,
उसका $A$ जो भी हो: \emph{थोलिआइटी}, फ़ेनर प्रवृत्ति का एक फ़ेरोबेसाल्ट। एंडेसाइट के
लिए, $A = 33.8$ पर रेखा (49.6, 35.0) और (52.5, 30.0) के बीच है, $F =
49.6 + 2.9\times1.2/5 = 50.3$ पर; $40.8 < 50.3$: \emph{कैल्क-क्षारीय}। (ग) मैग्नेटाइट,
जो ऑक्सीकृत \omterm{def:b2:magma-emplacement:magma}{मैग्मा} से आरंभ से ही फ़ेरिक लोहा बाहर ले जाता है और $F/M$ को
थोलिआइट जितनी तेज़ी से बढ़ने नहीं देता।
\end{solution}

\begin{solution}{exo:b2:magmatic-series:6}
MORB: $\mathrm{Nb_N} = 4/0.658 = 6.08$, $\mathrm{La_N} = 4.12/0.648 = 6.36$, अनुपात
0.956। महासागरीय-द्वीप बेसाल्ट: $38/0.658 = 57.8$, $30.8/0.648 = 47.5$, अनुपात 1.215। चाप
बेसाल्ट: $3/0.658 = 4.56$, $6.1/0.648 = 9.41$, अनुपात 0.484। केवल चाप बेसाल्ट में
नायोबियम विसंगति है, स्पाइडर आरेख पर Nb पर एक गर्त: नायोबियम अपने
पड़ोसी La की तुलना में क्षीण है: चाप \omterm{def:b2:magma-emplacement:magma}{मैग्मा} को स्लैब से तत्त्व मिले, पर नायोबियम
बहुत कम। MORB लगभग सपाट है; महासागरीय-द्वीप बेसाल्ट La की अपेक्षा
Nb में अधिक समृद्ध है।
\end{solution}

\begin{solution}{exo:b2:magmatic-series:7}
(क) $\mathrm{Ti} = 1.6\times0.5993\times10^4 = \qty{9589}{ppm}$; Ti/100 = 95.9, Zr
110, $3\mathrm Y = 96$; योग 301.9; निर्देशांक 31.8, 36.4, 31.8: क्षेत्र B। (ख) Ti
16\,780, अतः 167.8, 220, 87; योग 474.8; 35.3, 46.3, 18.3: क्षेत्र D। (ग) Ti 4794, अतः
47.9, 50, 60; योग 157.9; 30.3, 31.7, 38.0: क्षेत्र B। बेसाल्ट (ख) एक अंतःप्लेट बेसाल्ट है
(किसी महासागरीय द्वीप या महाद्वीप का); (क) और (ग), क्षेत्र B में, MORB, चाप
थोलिआइट या कैल्क-क्षारीय बेसाल्ट हो सकते हैं: आरेख निर्णय नहीं कर सकता, और कोई दूसरा आरेख (एक
स्पाइडर आरेख, Nb विसंगति) चाहिए।
\end{solution}

\begin{solution}{exo:b2:magmatic-series:8}
(क) MORB $142/164 = 0.866$, महासागरीय-द्वीप बेसाल्ट $4/173 = 0.023$, चाप बेसाल्ट $101/201 =
0.502$। (ख) समान पूर्व-प्रायिकताओं के साथ बेज़ नियम देता है $0.866/(0.866 + 0.023 + 0.502) =
0.866/1.391 = 0.623$। (ग) यह कि बेसाल्ट के महासागरीय-द्वीप बेसाल्ट होने की संभावना बहुत कम है
(प्रायिकता $0.023/1.391 = 0.017$), पर यह नहीं कि वह MORB है या चाप बेसाल्ट:
मोटे तौर पर पाँच में तीन संभावनाएँ, दो के विरुद्ध।
\end{solution}

\begin{solution}{exo:b2:magmatic-series:9}
प्रति \qty{100}{g} मोल: \ce{Na2O} $5.5/61.98 = 0.08874$, \ce{K2O} $4.5/94.20 =
0.04777$, \ce{Al2O3} $10.5/101.96 = 0.10298$। (क) $(0.08874 + 0.04777)/0.10298 =
0.13651/0.10298 = 1.326 > 1$: रायोलाइट \emph{अतिक्षारीय} है। (ख) फ़ेल्डस्पार
$0.10298$ मोल ले सकता है, क्षार ऑक्साइड के कुल $0.13651$ मोलों में से; आधिक्य, $0.03353$,
$0.03353/0.13651 = 0.246$ है, यानी क्षारों का एक चौथाई, जिसे अन्य खनिजों में जाना पड़ता है।
\end{solution}

\begin{solution}{exo:b2:magmatic-series:10}
(क) $s = (\lambda_F - \lambda_M)/(\lambda_F - \lambda_A)$: $\lambda_F = 0.5$ के लिए $(0.5
- 0.8)/(0.5 - 0.4) = -3.0$; $0.7$ के लिए $(-0.1)/(0.3) = -0.333$। (ख)
$s(\lambda_F - 0.4) = \lambda_F - 0.8$ को हल करने पर $\lambda_F = (0.4s - 0.8)/(s - 1)$ मिलता है:
आइसलैंड, $s = -1.10$: $(-0.44 - 0.8)/(-2.10) = 0.590$; कैस्केड, $s = -0.72$:
$(-0.288 - 0.8)/(-1.72) = 0.633$। (ग) कैस्केड का \omterm{def:b2:magma-emplacement:magma}{मैग्मा}, बड़े
$\lambda_F$ के साथ, अपना लोहा अधिक तेज़ी से खोता है: मैग्नेटाइट उसे बाहर ले जाता है। (घ) $\ln(F/M)$
दर $\lambda_M - \lambda_F$ से बढ़ता है और $\ln(A/F)$ दर $\lambda_F - \lambda_A$ से; अतः
$\ln(A/F)$ द्वारा प्राप्त हर इकाई पर $\ln(F/M)$ को $-s$ मिलता है: आइसलैंड में 1.10,
कैस्केड में 0.72। $F/M$ दोनों प्रवृत्तियों के अनुदिश बढ़ता है, पर थोलिआइटी प्रवृत्ति, यानी फ़ेनर
प्रवृत्ति, में अधिक तेज़ी से।
\end{solution}

\begin{solution}{exo:b2:magmatic-series:11}
(क) BHVO-2 के मोल, प्रति \qty{100}{g}, ये हैं: \ce{SiO2} 0.82557, \ce{Al2O3} 0.13182, \ce{FeO}
0.15521, \ce{MgO} 0.18007, \ce{CaO} 0.20328, \ce{Na2O} 0.03580, \ce{K2O} 0.00545।
ऑर्थोक्लेज़ 0.00545 और एल्बाइट 0.03580; शेष \ce{Al2O3} $0.13182 - 0.04125 = 0.09057$,
सब उतने ही \ce{CaO} के साथ एनॉर्थाइट में; शेष \ce{CaO} $0.11271$ डायोप्साइड में;
शेष $(\ce{Mg,Fe})\ce{O}$ $0.33528 - 0.11271 = 0.22257$ हाइपरस्थीन में। आवश्यक
सिलिका: $6\times0.00545 + 6\times0.03580 + 2\times0.09057 + 2\times0.11271 + 0.22257 =
0.87663$; उपलब्ध 0.82557; कमी 0.05106। हाइपरस्थीन से ओलिवीन में बदला गया हर $(\ce{Mg,Fe})\ce{O}$
आधा सिलिका बचाता है, अधिकतम $0.22257/2 = 0.11129$: उसमें से
$0.10212$ को बदलना पर्याप्त है, और $0.12045$ हाइपरस्थीन में रहता है। नॉर्म में ओलिवीन
और हाइपरस्थीन हैं, और न क्वार्ट्ज़ न नेफ़लीन: परिभाषा के अर्थ में चट्टान सिलिका-संतृप्त है,
एक हाइपरस्थीन-नॉर्मेटिव \emph{ओलिवीन थोलिआइट}, जो
\omterm{def:b2:magmatic-series:series}{थोलिआइटी श्रेणी} का है, जैसा हवाई ज्वालामुखी की ढाल-अवस्था उगलती है। (ख) मोल: 0.75732, 0.14712, 0.16008, 0.21092, 0.17832, 0.06131,
0.01592। ऑर्थोक्लेज़ 0.01592, एल्बाइट 0.06131, एनॉर्थाइट 0.06988, डायोप्साइड 0.10843, और हाइपरस्थीन के लिए $0.37100 - 0.10843 =
0.26257$; आवश्यक सिलिका $0.09552 + 0.36786 + 0.13976 + 0.21686 + 0.26257 = 1.08257$;
उपलब्ध 0.75732; कमी 0.32525। सारे हाइपरस्थीन को ओलिवीन में बदलने से
0.13129 बचता है, और 0.19396 शेष रहता है; एल्बाइट से नेफ़लीन में बदला गया \ce{Na2O} का हर मोल
4 \ce{SiO2} मुक्त करता है, अतः एल्बाइट के \ce{Na2O} के 0.06131 मोलों में से
$0.19396/4 = 0.04849$ नेफ़लीन बन जाते हैं। नॉर्म में नेफ़लीन: चट्टान
सिलिका-असंतृप्त है, एक \emph{क्षारीय बेसाल्ट}, जो \omterm{def:b2:magmatic-series:series}{क्षारीय श्रेणी} का है।
\end{solution}

\begin{solution}{exo:b2:magmatic-series:12}
(क) चाप बेसाल्ट $116/141 = 0.823$, MORB $21/101 = 0.208$, महासागरीय-द्वीप बेसाल्ट
$3/137 = 0.022$। (ख) $0.823/(0.823 + 0.208 + 0.022) = 0.823/1.053 = 0.782$।
(ग) अंश $0.1\times0.823 = 0.0823$; हर $0.0823 + 0.6\times0.208 +
0.3\times0.022 = 0.0823 + 0.1248 + 0.0066 = 0.2137$; प्रायिकता $0.385$। (घ)
एक ही चट्टान पर वही परीक्षण 0.78 या 0.39 देता है, इस पर निर्भर करते हुए कि जहाँ चट्टान मिली वहाँ
हर परिवेश कितना सामान्य है: \omterm{def:b2:magmatic-series:discrimination}{पृथक्करण आरेख} एक संभाविता देता है, और
भूवैज्ञानिक संदर्भ को शेष पूरा करना होता है।
\end{solution}

\begin{solution}{pb:b2:magmatic-series:1}
\noindent\textbf{भाग I।}

\textbf{1.} ऑक्साइड: लावा 1 99.17, लावा 2 98.65, लावा 3 98.20; वाष्पशीलों सहित
99.97, 99.80 और 99.45, सभी 100 के निकट। जल (0.70, 1.10, 1.20)
\qty{2}{\percent} से नीचे है और कार्बन डाइऑक्साइड (0.10, 0.05, 0.05) \qty{0.5}{\percent} से नीचे: तीनों
स्वीकार्य हैं।

\textbf{2.} गुणक $100/99.17 = 1.00837$: \ce{SiO2} 50.82, \ce{Na2O} 2.92, \ce{K2O}
0.63, क्षार 3.55।

\textbf{3.} लावा 2: $100/98.65 = 1.01368$, \ce{SiO2} 59.40, \ce{Na2O} 3.95, \ce{K2O}
1.17, क्षार 5.12। लावा 3: $100/98.20 = 1.01833$, \ce{SiO2} 66.50, \ce{Na2O} 4.48,
\ce{K2O} 1.88, क्षार 6.36।

\textbf{4.} लावा 1 (50.82, 3.55): बेसाल्ट। लावा 2 (59.40, 5.12; वहाँ O2 की ऊपरी सीमा
$5.9 + 1.1\times2.40/6 = 6.34$ पर है): एंडेसाइट। लावा 3 (66.50, 6.36; 63
और (69, 8) से (77, 0) तक की रेखा के बीच, (63, 7) से (69, 8) तक की रेखा से नीचे): डेसाइट।

\textbf{5.} कच्चे मान: (50.4, 3.52), (58.6, 5.05), (65.3, 6.25): वही क्षेत्र।
पुनर्परिकलन हर बिंदु को अधिकतम \qty{1.2}{\percent} सिलिका खिसकाता है और
यहाँ कुछ नहीं बदलता, \cref{exo:b2:magmatic-series:2} के विपरीत।

\textbf{6.} \omterm{def:b2:magmatic-series:qapf}{QAPF आरेख}, मापे गए मोड से। एंडेसाइटी संघटन वाली पातालीय चट्टान
प्लैजियोक्लेज़ में समृद्ध होती है, जिसमें क्षार-फ़ेल्डस्पार या क्वार्ट्ज़ कम होता है: उसे
क्षेत्र 10 या 10* में पड़ना चाहिए, एक डायोराइट या क्वार्ट्ज़ डायोराइट (उसका मोड तय करता है)।

\noindent\textbf{भाग II।}

\textbf{7.} लावा 1: रेखाएँ 0.43 और 1.47 पर, $\ce{K2O} = 0.63$: मध्यम-K। लावा 2: 0.81
और 2.28, 1.17: मध्यम-K। लावा 3: 1.13 और 2.96, 1.88: मध्यम-K।

\textbf{8.} वाष्पशील-मुक्त \ce{FeO_T} और \ce{MgO}: 8.67 और 7.97; 6.18 और 3.85; 4.07
और 1.63। लावा 1: $3.55 + 8.67 + 7.97 = 20.19$; $A = 17.6$, $F = 42.9$, $M = 39.5$।
लावा 2: 15.15; 33.8, 40.8, 25.4। लावा 3: 12.06; 52.7, 33.7, 13.5।

\textbf{9.} लावा 1, $A = 17.6$: (42.8, 17.1) और (46.9, 18.1) के बीच रेखा
$F = 42.8 + 4.1\times0.5/1.0 = 44.9$ पर है; $42.9 < 44.9$, कैल्क-क्षारीय ओर, रेखा के
निकट। लावा 2, $A = 33.8$: $F = 49.6 + 2.9\times1.2/5 = 50.3$; $40.8 < 50.3$,
कैल्क-क्षारीय। लावा 3, $A = 52.7$: $F = 40.8 - 2.9\times2.7/5 = 39.2$; $33.7 < 39.2$,
कैल्क-क्षारीय।

\textbf{10.} $\ce{FeO_T}/\ce{MgO}$: $8.67/7.97 = 1.09$, $6.18/3.85 = 1.61$, $4.07/1.63 =
2.50$। अनुपात बढ़ता है, जैसा हर विभेदित होते बेसाल्ट में बढ़ता है, पर लोहा स्वयं
8.67 से 4.07 तक घटता है, और त्रिभुज पर $F$ 42.9 से 33.7 तक घटता है: कोई लौह-संवर्धन
नहीं। थोलिआइटी समूह पहले F की ओर चढ़ता, जैसा
\cref{exo:b2:magmatic-series:5} का फ़ेरोबेसाल्ट करता है।

\textbf{11.} \omterm{def:b2:magmatic-series:series}{कैल्क-क्षारीय श्रेणी}। ऑक्सीकृत \omterm{def:b2:magma-emplacement:magma}{मैग्मा} से क्रिस्टलित होते मैग्नेटाइट ने
मैग्नीशियम के साथ लोहा हटाया; ऑक्सीकृत स्पीशीज़ किसी क्षेपण क्षेत्र के निर्जलित होते स्लैब से
आती हैं।

\noindent\textbf{भाग III।}

\textbf{12.} $\mathrm{Nb_N} = 2.0/0.658 = 3.04$, $\mathrm{La_N} = 7/0.648 = 10.80$,
$\mathrm{Th_N} = 1.1/0.0795 = 13.84$।

\textbf{13.} $3.04/10.80 = 0.281$, MORB (0.956) और
महासागरीय-द्वीप बेसाल्टों (1.215) की माध्यिकाओं से बहुत नीचे, और माध्यिका चाप बेसाल्ट (0.484) से भी नीचे: एक गहरा
नायोबियम गर्त।

\textbf{14.} $13.84/3.04 = 4.55$; स्पाइडर आरेख के माध्यिका MORB का $2.64/6.08 =
0.43$ है। थोरियम को, चाप बेसाल्टों में समृद्ध अन्य तत्त्वों की तरह,
धँसते स्लैब ने जोड़ा।

\textbf{15.} वाष्पशील-मुक्त \ce{TiO2} $0.95\times1.00837 = 0.958$; $\mathrm{Ti} =
0.958\times5993 = \qty{5741}{ppm}$; Ti/100 = 57.4, Zr 70, $3\mathrm Y = 54$, योग 181.4:
31.6, 38.6, 29.8, क्षेत्र B में, जो MORB, चाप थोलिआइटों और कैल्क-क्षारीय बेसाल्टों में साझा है।
यह अंतःप्लेट बेसाल्ट को बाहर करता है, पर कटक और चाप के बीच निर्णय नहीं करता।
आरेख बेसाल्टों के लिए बनाया गया था: एंडेसाइट और डेसाइट में टाइटेनियम ऑक्साइडों द्वारा
हटाया जा चुका है और ज़िरकोनियम क्रिस्टलन से सांद्रित हुआ है, और उनके क्षेत्र (लावा 2
C में पड़ता, लावा 3 हर क्षेत्र के बाहर) का कोई अर्थ नहीं होता।

\textbf{16.} एक कैल्क-क्षारीय, मध्यम-K श्रेणी, गहरे नायोबियम गर्त और ऊँचे
$\mathrm{Th_N}/\mathrm{Nb_N}$ के साथ, ऐसे ज्वालामुखी पर जिसने बेसाल्ट से डेसाइट तक उगला: ज्वालामुखी
किसी क्षेपण क्षेत्र के ऊपर एक ज्वालामुखी चाप में बढ़ा।

\noindent\textbf{भाग IV।}

\textbf{17.} पूर्णतः \omterm{def:b2:elements-in-the-earth:partition}{असंगत तत्त्व} के लिए $C_L = C_0/F$: $F = 0.63/1.17 =
0.538$।

\textbf{18.} रुबिडियम $10/20 = 0.500$; थोरियम $1.1/2.1 = 0.524$। (\omterm{def:b2:elements-in-the-earth:trace}{अनुरेखी
तत्त्व}, जो विश्लेषित रूप में दिए गए हैं, पुनर्परिकलन से
\qty{1}{\percent} से कम बदलते।)

\textbf{19.} $C_L/C_0 = F^{D - 1}$: $125/70 = 1.786 = F^{-0.9}$, $F = 1.786^{-1/0.9} =
0.525$।

\textbf{20.} पोटाश से, लावा 2 से लावा 3 तक $F = 1.17/1.88 = 0.622$। स्ट्रॉन्शियम:
$380/480 = 0.792 = 0.622^{D - 1}$, अतः $D - 1 = \ln 0.792/\ln 0.622 = -0.234/-0.474 =
0.49$ और $D \approx 1.49$: स्ट्रॉन्शियम संगत हो गया है। प्लैजियोक्लेज़, जो
स्ट्रॉन्शियम ग्रहण करता है, प्रचुरता से क्रिस्टलित हो रहा था।

\textbf{21.} पोटाश से, $F = 0.63/1.88 = 0.335$ (रुबिडियम $10/32 = 0.31$, थोरियम
$1.1/3.4 = 0.32$): डेसाइट से पहले बेसाल्ट का लगभग दो-तिहाई भाग, $1 - 0.335 = 0.665$,
क्रिस्टलित हुआ।

\textbf{22.} आत्मसात पोटाश लावा 2 में \ce{K2O} को अकेले क्रिस्टलन से मिलने वाले मान से
ऊपर उठाता, अतः $F = 0.63/1.17$ बहुत कम आता, और क्रिस्टलों का
द्रव्यमान बहुत अधिक।

\textbf{23.} $(0.538 + 0.500 + 0.524 + 0.525)/4 = 2.087/4 = 0.522$; क्रिस्टल $1 - 0.522
= 0.478$ बेसाल्ट का, यानी एंडेसाइट बनाने के लिए \textbf{\qty[text-series-to-math=true]{48}{kg}
क्रिस्टल, बेसाल्ट के हर \qty{100}{kg} पर}।
\end{solution}
